\chapter{After the Screening}\label{ch:after-screening}

A film continues after it ends, in what its audience says about it. In ordinary cinema that conversation starts from a common object. People disagree about what a film meant, what a character intended, whether the ending worked, but they agree about what was shown. In selective cinema that agreement is gone, and the conversation changes character. This chapter examines what happens after a selective screening, and then turns to the protection of what one realization must keep from another's viewers.

\section{Which one did you see?}

After an ordinary film, ``What did you think happened?'' is a question about interpretation. After a selective one, it cannot be separated from a prior question: ``Which version did you see?'' Two people may disagree because they read the same evidence differently, or because they were shown different evidence. Until they know which, they cannot tell whether they are disagreeing at all.

The effect depends on the family. After a genre pair, the evidence is nearly the same and the conversation is about tone: one found it frightening, the other funny, and each learns what the other's film made of the same events. After a depth family, the conversation is about understanding, and the viewer of the expert realization may explain what the beginner realization spelled out, or be surprised by how much it did. After a family of endings, the conversation begins with incompatible outcomes, and the audience must reconstruct the family together from its parts. \Cref{ch:endings} called this differentiated memory: people who shared a room leave with memories of the same film that cannot both be true of one sequence.

\section{Three regimes of knowledge}

What the audience knows about the variation, before and after, shapes all of this. Three regimes are possible.

In the first, selections are \emph{declared}. Everyone knows the family exists, what its realizations are, and who chose which; the dials are visible, or the choice is announced. The conversation afterward is informed from its first sentence.

In the second, selections are \emph{private}. Everyone knows the family exists and what its realizations are, but each viewer's choice is their own unless they share it. The conversation afterward includes the choice of whether to say.

In the third, the variation is \emph{undisclosed}. Viewers do not know that others are seeing something different. They discover it, if they do, only in conversation, when a companion describes a scene they do not recognize.

The first two regimes are legitimate and offer different experiences. Declared selections make the family a public game; private selections protect the viewer's choice, which may say something about them, as \Cref{ch:depth} noted of expertise. The third regime is different in kind. It turns a work about difference into an instrument of deception: people believe they attended one event and saw one film, and they did not. Some artistic uses might want the shock of discovery, and a work could reasonably withhold which realization a viewer will receive. But the existence of variation should never be concealed. A viewer is entitled to know that the people beside them may be seeing something else. \Cref{ch:persuasion} develops this into a disclosure requirement.

\section{The conversation as final act}

Under the declared and private regimes, the conversation afterward is not an accident of the form. It can be designed for. A family whose realizations each reveal part of a larger picture invites its audience to assemble it. A mystery whose realizations follow different investigators can leave each viewer with a piece of the solution that only the group can complete. A family of endings can be written so that the endings, compared, reveal something none shows alone.

The strongest uses of selective cinema may treat the conversation after the screening as the work's final act: a part performed by the audience, with material the film distributed among them. That is not possible in any form where everyone receives the same thing, and it is only awkwardly possible in forms where people experience a work separately and at different times. It requires exactly what selective cinema provides: people who were together, on one clock, in one story, and were shown different parts of it.

\section{Protected variation}

The conversation depends on something the apparatus cannot fully guarantee: that each viewer actually saw only their own realization. \Cref{ch:phase} showed that leakage between streams can be reduced but not eliminated, and that it has two measures, optical and informational. This section turns that observation into a standard.

Not all variation needs protection. A genre pair whose realizations show the same events with different performances loses nothing if a viewer glimpses the other performance for a fraction of a second. Call such variation \emph{casual}. A depth family whose expert realization shows a formula the beginner realization explains later loses a little if the formula leaks early; call it \emph{sensitive}. A family whose realizations withhold a revelation from each other, the identity of a culprit, an ending, a death, loses everything if the revelation leaks; call it \emph{protected}.

The harm of a leak is roughly the product of two factors: how much the leaked information matters, and how recognizable the leak is. A faint luminance ghost carrying no recognizable content is harmless in any realization. A clear glimpse of a face is harmless in a casual family and ruinous if the face is the protected revelation. A standard for leakage should therefore be set by narrative consequence, not by optics alone. Casual variation tolerates the leakage of ordinary shutter glasses. Sensitive variation needs leakage low enough that no leaked frame is legible. Protected variation needs leakage below legibility, protection of the composite, and particular attention to the moments immediately before and after the revelation.

The protection profile that \Cref{ch:editing} made part of the editor's deliverables is the implementation. It tells the exhibition system when to lengthen blanking intervals, reduce the number of active streams, or suppress the composite, so that the strict standard applies only where it is needed and its costs in light are paid only there.

\section{What protection cannot do}

Protection by timing has the limits already identified in \Cref{ch:unfiltered}. It protects against the eye and not against cameras. A viewer determined to see another realization can adjust their glasses, borrow a companion's, or film the screen. Protected variation in a selective screening is protection against accident, not against intent.

That is enough for its purpose. The point of protecting a revelation is to let viewers who chose not to see it be surprised by it, or never meet it at all. It is not to prevent anyone from ever finding out. The conversation afterward will reveal it in any case, if the viewers want it revealed. Protection keeps the revelation for the people who chose it, at the moment the film intended, and leaves what happens next to the audience. A system that promised more, that a realization could be kept from people who wanted to see it, would be offering concealment, and \Cref{ch:perceptual-rights} examines what that requires and why it changes what the glasses are.
